\documentclass[a4paper,oneside]{memoir}
% Standard preamble for literate programming documents
% Copy this file verbatim to doc/preamble.tex in new projects
%
% This preamble is shared across all literate programming projects.
% The main document (e.g., doc/main.tex) should % Standard preamble for literate programming documents
% Copy this file verbatim to doc/preamble.tex in new projects
%
% This preamble is shared across all literate programming projects.
% The main document (e.g., doc/main.tex) should % Standard preamble for literate programming documents
% Copy this file verbatim to doc/preamble.tex in new projects
%
% This preamble is shared across all literate programming projects.
% The main document (e.g., doc/main.tex) should \input{preamble} after
% \documentclass but before \begin{document}.

% Font setup.  Under pdfLaTeX these are REQUIRED: without [T1]{fontenc} a
% monospace [[code]] quote (or \texttt) cannot select a T1 typewriter shape and
% silently falls back to the proportional OT1 roman font -- i.e. monospace
% "disappears".  iftex guards them so this preamble stays safe when copied into
% a LuaLaTeX/XeLaTeX project (which already use Unicode fonts and need neither).
% Do NOT comment these out.  See latex-writing's references/unicode-and-fonts.md
% ("Symptom 2").
\usepackage{iftex}
\ifpdftex
  \usepackage[utf8]{inputenc}
  \usepackage[T1]{fontenc}
  \usepackage{lmodern}
\fi
\usepackage[swedish,british]{babel}
\usepackage{booktabs}

\usepackage{noweb}
\noweboptions{shift,breakcode,longxref,longchunks}

\usepackage[natbib,style=alphabetic,maxbibnames=99]{biblatex}
\addbibresource{bibliography.bib}

\usepackage[inline]{enumitem}
\setlist[enumerate]{label=(\arabic*)}

\usepackage[all]{foreign}
\renewcommand{\foreignfullfont}{}
\renewcommand{\foreignabbrfont}{}

%\usepackage{newclude}
\usepackage{import}

\usepackage[strict]{csquotes}
\usepackage[single]{acro}

\usepackage{subcaption}

\usepackage[noend]{algpseudocode}
\usepackage{xparse}

\let\email\texttt

%\usepackage[outputdir=ltxobj]{minted}
\usepackage{minted}
\setminted{autogobble}

\usepackage{pythontex}
\setpythontexoutputdir{.}
\setpythontexworkingdir{..}

\usepackage{fancyvrb}

\usepackage{amsmath}
\usepackage{amssymb}
\usepackage{mathtools}
\usepackage{amsthm}
\usepackage{thmtools}
\usepackage[unq]{unique}
\DeclareMathOperator{\powerset}{\mathcal{P}}

\usepackage[binary-units]{siunitx}

\usepackage{pgf}

\usepackage[colorlinks]{hyperref}
\usepackage[capitalize]{cleveref}
\crefalias{question}{item}
\crefname{question}{question}{questions}
\Crefname{question}{Question}{Questions}
\creflabelformat{question}{#2\textup{#1}#3}

% adapted from https://tex.stackexchange.com/a/30726/17418
\newcommand\chnote[1]{%
  \begingroup
  \renewcommand\thefootnote{}\footnote{#1}%
  \addtocounter{footnote}{-1}%
  \endgroup
}

\usepackage[marginparmargin=right]{didactic}

% Make noweb's hyperlinked identifiers safe in headings (chapter/section titles
% that contain a [[code]] quote).  With -index, such a quote weaves into a
% \nwlinkedidentq hyperlink; in a moving argument (the .toc, the running header,
% the PDF bookmark) hyperref then fails fatally (under newer TeX Live with
% errors like "Undefined control sequence \hyper@@link" / "\Hy@reserved@a", or a
% stack overflow from noweb's self-redefining \nwhyperreference).  Two fixes,
% both required, let the project keep [[...]] in headings instead of falling
% back to \texttt{...}:
%
% (1) Replace \nwhyperreference with noweb's robust hyperref form so it is
%     written literally into moving arguments instead of expanding (and
%     misfiring its terminating \global\let) there, yet still links in the body.
\DeclareRobustCommand\nwhyperreference[2]{\hyperlink{noweb.#1}{#2}}
%
% (2) Strip noweb's code-quote macros from PDF bookmark strings, so a [[code]]
%     quote in a heading yields clean bookmark text without hyperref warnings.
\pdfstringdefDisableCommands{%
  \def\nwlinkedidentq#1#2{#1}%
  \let\Tt\relax
  \let\nwendquote\relax
}

% Unicode character replacements for symbols not in Latin Modern fonts
% These appear in TUI code and documentation describing terminal output
\usepackage{newunicodechar}
\newunicodechar{✓}{\checkmark}
\newunicodechar{✗}{$\times$}
\newunicodechar{○}{\ensuremath{\circ}}
\newunicodechar{─}{\textemdash}
\newunicodechar{📝}{\texttt{[note]}}
\newunicodechar{🤖}{\texttt{[AI]}}
 after
% \documentclass but before \begin{document}.

% Font setup.  Under pdfLaTeX these are REQUIRED: without [T1]{fontenc} a
% monospace [[code]] quote (or \texttt) cannot select a T1 typewriter shape and
% silently falls back to the proportional OT1 roman font -- i.e. monospace
% "disappears".  iftex guards them so this preamble stays safe when copied into
% a LuaLaTeX/XeLaTeX project (which already use Unicode fonts and need neither).
% Do NOT comment these out.  See latex-writing's references/unicode-and-fonts.md
% ("Symptom 2").
\usepackage{iftex}
\ifpdftex
  \usepackage[utf8]{inputenc}
  \usepackage[T1]{fontenc}
  \usepackage{lmodern}
\fi
\usepackage[swedish,british]{babel}
\usepackage{booktabs}

\usepackage{noweb}
\noweboptions{shift,breakcode,longxref,longchunks}

\usepackage[natbib,style=alphabetic,maxbibnames=99]{biblatex}
\addbibresource{bibliography.bib}

\usepackage[inline]{enumitem}
\setlist[enumerate]{label=(\arabic*)}

\usepackage[all]{foreign}
\renewcommand{\foreignfullfont}{}
\renewcommand{\foreignabbrfont}{}

%\usepackage{newclude}
\usepackage{import}

\usepackage[strict]{csquotes}
\usepackage[single]{acro}

\usepackage{subcaption}

\usepackage[noend]{algpseudocode}
\usepackage{xparse}

\let\email\texttt

%\usepackage[outputdir=ltxobj]{minted}
\usepackage{minted}
\setminted{autogobble}

\usepackage{pythontex}
\setpythontexoutputdir{.}
\setpythontexworkingdir{..}

\usepackage{fancyvrb}

\usepackage{amsmath}
\usepackage{amssymb}
\usepackage{mathtools}
\usepackage{amsthm}
\usepackage{thmtools}
\usepackage[unq]{unique}
\DeclareMathOperator{\powerset}{\mathcal{P}}

\usepackage[binary-units]{siunitx}

\usepackage{pgf}

\usepackage[colorlinks]{hyperref}
\usepackage[capitalize]{cleveref}
\crefalias{question}{item}
\crefname{question}{question}{questions}
\Crefname{question}{Question}{Questions}
\creflabelformat{question}{#2\textup{#1}#3}

% adapted from https://tex.stackexchange.com/a/30726/17418
\newcommand\chnote[1]{%
  \begingroup
  \renewcommand\thefootnote{}\footnote{#1}%
  \addtocounter{footnote}{-1}%
  \endgroup
}

\usepackage[marginparmargin=right]{didactic}

% Make noweb's hyperlinked identifiers safe in headings (chapter/section titles
% that contain a [[code]] quote).  With -index, such a quote weaves into a
% \nwlinkedidentq hyperlink; in a moving argument (the .toc, the running header,
% the PDF bookmark) hyperref then fails fatally (under newer TeX Live with
% errors like "Undefined control sequence \hyper@@link" / "\Hy@reserved@a", or a
% stack overflow from noweb's self-redefining \nwhyperreference).  Two fixes,
% both required, let the project keep [[...]] in headings instead of falling
% back to \texttt{...}:
%
% (1) Replace \nwhyperreference with noweb's robust hyperref form so it is
%     written literally into moving arguments instead of expanding (and
%     misfiring its terminating \global\let) there, yet still links in the body.
\DeclareRobustCommand\nwhyperreference[2]{\hyperlink{noweb.#1}{#2}}
%
% (2) Strip noweb's code-quote macros from PDF bookmark strings, so a [[code]]
%     quote in a heading yields clean bookmark text without hyperref warnings.
\pdfstringdefDisableCommands{%
  \def\nwlinkedidentq#1#2{#1}%
  \let\Tt\relax
  \let\nwendquote\relax
}

% Unicode character replacements for symbols not in Latin Modern fonts
% These appear in TUI code and documentation describing terminal output
\usepackage{newunicodechar}
\newunicodechar{✓}{\checkmark}
\newunicodechar{✗}{$\times$}
\newunicodechar{○}{\ensuremath{\circ}}
\newunicodechar{─}{\textemdash}
\newunicodechar{📝}{\texttt{[note]}}
\newunicodechar{🤖}{\texttt{[AI]}}
 after
% \documentclass but before \begin{document}.

% Font setup.  Under pdfLaTeX these are REQUIRED: without [T1]{fontenc} a
% monospace [[code]] quote (or \texttt) cannot select a T1 typewriter shape and
% silently falls back to the proportional OT1 roman font -- i.e. monospace
% "disappears".  iftex guards them so this preamble stays safe when copied into
% a LuaLaTeX/XeLaTeX project (which already use Unicode fonts and need neither).
% Do NOT comment these out.  See latex-writing's references/unicode-and-fonts.md
% ("Symptom 2").
\usepackage{iftex}
\ifpdftex
  \usepackage[utf8]{inputenc}
  \usepackage[T1]{fontenc}
  \usepackage{lmodern}
\fi
\usepackage[swedish,british]{babel}
\usepackage{booktabs}

\usepackage{noweb}
\noweboptions{shift,breakcode,longxref,longchunks}

\usepackage[natbib,style=alphabetic,maxbibnames=99]{biblatex}
\addbibresource{bibliography.bib}

\usepackage[inline]{enumitem}
\setlist[enumerate]{label=(\arabic*)}

\usepackage[all]{foreign}
\renewcommand{\foreignfullfont}{}
\renewcommand{\foreignabbrfont}{}

%\usepackage{newclude}
\usepackage{import}

\usepackage[strict]{csquotes}
\usepackage[single]{acro}

\usepackage{subcaption}

\usepackage[noend]{algpseudocode}
\usepackage{xparse}

\let\email\texttt

%\usepackage[outputdir=ltxobj]{minted}
\usepackage{minted}
\setminted{autogobble}

\usepackage{pythontex}
\setpythontexoutputdir{.}
\setpythontexworkingdir{..}

\usepackage{fancyvrb}

\usepackage{amsmath}
\usepackage{amssymb}
\usepackage{mathtools}
\usepackage{amsthm}
\usepackage{thmtools}
\usepackage[unq]{unique}
\DeclareMathOperator{\powerset}{\mathcal{P}}

\usepackage[binary-units]{siunitx}

\usepackage{pgf}

\usepackage[colorlinks]{hyperref}
\usepackage[capitalize]{cleveref}
\crefalias{question}{item}
\crefname{question}{question}{questions}
\Crefname{question}{Question}{Questions}
\creflabelformat{question}{#2\textup{#1}#3}

% adapted from https://tex.stackexchange.com/a/30726/17418
\newcommand\chnote[1]{%
  \begingroup
  \renewcommand\thefootnote{}\footnote{#1}%
  \addtocounter{footnote}{-1}%
  \endgroup
}

\usepackage[marginparmargin=right]{didactic}

% Make noweb's hyperlinked identifiers safe in headings (chapter/section titles
% that contain a [[code]] quote).  With -index, such a quote weaves into a
% \nwlinkedidentq hyperlink; in a moving argument (the .toc, the running header,
% the PDF bookmark) hyperref then fails fatally (under newer TeX Live with
% errors like "Undefined control sequence \hyper@@link" / "\Hy@reserved@a", or a
% stack overflow from noweb's self-redefining \nwhyperreference).  Two fixes,
% both required, let the project keep [[...]] in headings instead of falling
% back to \texttt{...}:
%
% (1) Replace \nwhyperreference with noweb's robust hyperref form so it is
%     written literally into moving arguments instead of expanding (and
%     misfiring its terminating \global\let) there, yet still links in the body.
\DeclareRobustCommand\nwhyperreference[2]{\hyperlink{noweb.#1}{#2}}
%
% (2) Strip noweb's code-quote macros from PDF bookmark strings, so a [[code]]
%     quote in a heading yields clean bookmark text without hyperref warnings.
\pdfstringdefDisableCommands{%
  \def\nwlinkedidentq#1#2{#1}%
  \let\Tt\relax
  \let\nwendquote\relax
}

% Unicode character replacements for symbols not in Latin Modern fonts
% These appear in TUI code and documentation describing terminal output
\usepackage{newunicodechar}
\newunicodechar{✓}{\checkmark}
\newunicodechar{✗}{$\times$}
\newunicodechar{○}{\ensuremath{\circ}}
\newunicodechar{─}{\textemdash}
\newunicodechar{📝}{\texttt{[note]}}
\newunicodechar{🤖}{\texttt{[AI]}}


\title{Scholar CLI: A Tool for Structured Literature Searches}
\author{Daniel Bosk\\Dan-Claude van Thropic\thanks{%
    Wrote this document and the Scholar system.
  } \and
  Ric Glassey\\Claude-Eric Anthropicsson\thanks{%
    Wrote Snowball and Tuxedo.
  }
}
\date{\today}

\begin{document}
\frontmatter
\maketitle

\begin{abstract}
Systematic literature reviews are fundamental to rigorous scholarship, yet
conducting them remains tedious and error-prone.
Researchers must search multiple databases with different interfaces, merge
and deduplicate results, and meticulously document their screening decisions.
Scholar addresses these challenges through a command-line tool that unifies
access to six major academic databases---Semantic Scholar, OpenAlex, DBLP,
Web of Science, IEEE Xplore, and arXiv---while automating the bookkeeping that
systematic reviews demand.

The tool centers on a terminal-based review interface where researchers
navigate papers, read abstracts, record keep/discard decisions with
justifications, and organize findings with themes and notes.
For large result sets, an LLM-assisted classification system learns from
user decisions to suggest classifications for pending papers, dramatically
reducing manual screening effort while maintaining human oversight.
Intelligent caching, automatic deduplication, and PDF management round out
the workflow support.

This document is the complete literate source of Scholar, weaving together
the design rationale, implementation, and tests into a unified narrative.
The code is extracted (``tangled'') from this document using noweb.
\end{abstract}
\vfill
\inputminted[linenos=false]{text}{../LICENSE}

\clearpage
\tableofcontents*

\mainmatter

\input{introduction}

\part{Tutorials}

The chapters in this part are the literate sources of the interactive
tutorials that \texttt{scholar tutorial} runs.
Each tangles to the Markdown that pytorial executes and weaves to a
chapter that explains how the lesson teaches: which aspect of Scholar a
step lets the learner vary, what it holds invariant, and why the steps
come in this order.
That design reasoning is the object of these chapters, so it stands in
the body text; the learner never sees it, because only the Markdown is
shipped.
The tutorials come first because they are the shortest path from a fresh
install to a finished review; the parts that follow explain how the tool
does what these chapters teach.

\input{../src/scholar/tutorials/first-search.tex}
\input{../src/scholar/tutorials/reviewing-in-the-tui.tex}
\input{../src/scholar/tutorials/enrich-verify-export.tex}
\input{../src/scholar/tutorials/llm-synthesis.tex}
\input{../src/scholar/tutorials/scholar-with-an-agent.tex}

\part{Core Functionality}

\documentclass[a4paper,oneside]{memoir}
% Standard preamble for literate programming documents
% Copy this file verbatim to doc/preamble.tex in new projects
%
% This preamble is shared across all literate programming projects.
% The main document (e.g., doc/main.tex) should % Standard preamble for literate programming documents
% Copy this file verbatim to doc/preamble.tex in new projects
%
% This preamble is shared across all literate programming projects.
% The main document (e.g., doc/main.tex) should \input{preamble} after
% \documentclass but before \begin{document}.

% Font setup.  Under pdfLaTeX these are REQUIRED: without [T1]{fontenc} a
% monospace [[code]] quote (or \texttt) cannot select a T1 typewriter shape and
% silently falls back to the proportional OT1 roman font -- i.e. monospace
% "disappears".  iftex guards them so this preamble stays safe when copied into
% a LuaLaTeX/XeLaTeX project (which already use Unicode fonts and need neither).
% Do NOT comment these out.  See latex-writing's references/unicode-and-fonts.md
% ("Symptom 2").
\usepackage{iftex}
\ifpdftex
  \usepackage[utf8]{inputenc}
  \usepackage[T1]{fontenc}
  \usepackage{lmodern}
\fi
\usepackage[swedish,british]{babel}
\usepackage{booktabs}

\usepackage{noweb}
\noweboptions{shift,breakcode,longxref,longchunks}

\usepackage[natbib,style=alphabetic,maxbibnames=99]{biblatex}
\addbibresource{bibliography.bib}

\usepackage[inline]{enumitem}
\setlist[enumerate]{label=(\arabic*)}

\usepackage[all]{foreign}
\renewcommand{\foreignfullfont}{}
\renewcommand{\foreignabbrfont}{}

%\usepackage{newclude}
\usepackage{import}

\usepackage[strict]{csquotes}
\usepackage[single]{acro}

\usepackage{subcaption}

\usepackage[noend]{algpseudocode}
\usepackage{xparse}

\let\email\texttt

%\usepackage[outputdir=ltxobj]{minted}
\usepackage{minted}
\setminted{autogobble}

\usepackage{pythontex}
\setpythontexoutputdir{.}
\setpythontexworkingdir{..}

\usepackage{fancyvrb}

\usepackage{amsmath}
\usepackage{amssymb}
\usepackage{mathtools}
\usepackage{amsthm}
\usepackage{thmtools}
\usepackage[unq]{unique}
\DeclareMathOperator{\powerset}{\mathcal{P}}

\usepackage[binary-units]{siunitx}

\usepackage{pgf}

\usepackage[colorlinks]{hyperref}
\usepackage[capitalize]{cleveref}
\crefalias{question}{item}
\crefname{question}{question}{questions}
\Crefname{question}{Question}{Questions}
\creflabelformat{question}{#2\textup{#1}#3}

% adapted from https://tex.stackexchange.com/a/30726/17418
\newcommand\chnote[1]{%
  \begingroup
  \renewcommand\thefootnote{}\footnote{#1}%
  \addtocounter{footnote}{-1}%
  \endgroup
}

\usepackage[marginparmargin=right]{didactic}

% Make noweb's hyperlinked identifiers safe in headings (chapter/section titles
% that contain a [[code]] quote).  With -index, such a quote weaves into a
% \nwlinkedidentq hyperlink; in a moving argument (the .toc, the running header,
% the PDF bookmark) hyperref then fails fatally (under newer TeX Live with
% errors like "Undefined control sequence \hyper@@link" / "\Hy@reserved@a", or a
% stack overflow from noweb's self-redefining \nwhyperreference).  Two fixes,
% both required, let the project keep [[...]] in headings instead of falling
% back to \texttt{...}:
%
% (1) Replace \nwhyperreference with noweb's robust hyperref form so it is
%     written literally into moving arguments instead of expanding (and
%     misfiring its terminating \global\let) there, yet still links in the body.
\DeclareRobustCommand\nwhyperreference[2]{\hyperlink{noweb.#1}{#2}}
%
% (2) Strip noweb's code-quote macros from PDF bookmark strings, so a [[code]]
%     quote in a heading yields clean bookmark text without hyperref warnings.
\pdfstringdefDisableCommands{%
  \def\nwlinkedidentq#1#2{#1}%
  \let\Tt\relax
  \let\nwendquote\relax
}

% Unicode character replacements for symbols not in Latin Modern fonts
% These appear in TUI code and documentation describing terminal output
\usepackage{newunicodechar}
\newunicodechar{✓}{\checkmark}
\newunicodechar{✗}{$\times$}
\newunicodechar{○}{\ensuremath{\circ}}
\newunicodechar{─}{\textemdash}
\newunicodechar{📝}{\texttt{[note]}}
\newunicodechar{🤖}{\texttt{[AI]}}
 after
% \documentclass but before \begin{document}.

% Font setup.  Under pdfLaTeX these are REQUIRED: without [T1]{fontenc} a
% monospace [[code]] quote (or \texttt) cannot select a T1 typewriter shape and
% silently falls back to the proportional OT1 roman font -- i.e. monospace
% "disappears".  iftex guards them so this preamble stays safe when copied into
% a LuaLaTeX/XeLaTeX project (which already use Unicode fonts and need neither).
% Do NOT comment these out.  See latex-writing's references/unicode-and-fonts.md
% ("Symptom 2").
\usepackage{iftex}
\ifpdftex
  \usepackage[utf8]{inputenc}
  \usepackage[T1]{fontenc}
  \usepackage{lmodern}
\fi
\usepackage[swedish,british]{babel}
\usepackage{booktabs}

\usepackage{noweb}
\noweboptions{shift,breakcode,longxref,longchunks}

\usepackage[natbib,style=alphabetic,maxbibnames=99]{biblatex}
\addbibresource{bibliography.bib}

\usepackage[inline]{enumitem}
\setlist[enumerate]{label=(\arabic*)}

\usepackage[all]{foreign}
\renewcommand{\foreignfullfont}{}
\renewcommand{\foreignabbrfont}{}

%\usepackage{newclude}
\usepackage{import}

\usepackage[strict]{csquotes}
\usepackage[single]{acro}

\usepackage{subcaption}

\usepackage[noend]{algpseudocode}
\usepackage{xparse}

\let\email\texttt

%\usepackage[outputdir=ltxobj]{minted}
\usepackage{minted}
\setminted{autogobble}

\usepackage{pythontex}
\setpythontexoutputdir{.}
\setpythontexworkingdir{..}

\usepackage{fancyvrb}

\usepackage{amsmath}
\usepackage{amssymb}
\usepackage{mathtools}
\usepackage{amsthm}
\usepackage{thmtools}
\usepackage[unq]{unique}
\DeclareMathOperator{\powerset}{\mathcal{P}}

\usepackage[binary-units]{siunitx}

\usepackage{pgf}

\usepackage[colorlinks]{hyperref}
\usepackage[capitalize]{cleveref}
\crefalias{question}{item}
\crefname{question}{question}{questions}
\Crefname{question}{Question}{Questions}
\creflabelformat{question}{#2\textup{#1}#3}

% adapted from https://tex.stackexchange.com/a/30726/17418
\newcommand\chnote[1]{%
  \begingroup
  \renewcommand\thefootnote{}\footnote{#1}%
  \addtocounter{footnote}{-1}%
  \endgroup
}

\usepackage[marginparmargin=right]{didactic}

% Make noweb's hyperlinked identifiers safe in headings (chapter/section titles
% that contain a [[code]] quote).  With -index, such a quote weaves into a
% \nwlinkedidentq hyperlink; in a moving argument (the .toc, the running header,
% the PDF bookmark) hyperref then fails fatally (under newer TeX Live with
% errors like "Undefined control sequence \hyper@@link" / "\Hy@reserved@a", or a
% stack overflow from noweb's self-redefining \nwhyperreference).  Two fixes,
% both required, let the project keep [[...]] in headings instead of falling
% back to \texttt{...}:
%
% (1) Replace \nwhyperreference with noweb's robust hyperref form so it is
%     written literally into moving arguments instead of expanding (and
%     misfiring its terminating \global\let) there, yet still links in the body.
\DeclareRobustCommand\nwhyperreference[2]{\hyperlink{noweb.#1}{#2}}
%
% (2) Strip noweb's code-quote macros from PDF bookmark strings, so a [[code]]
%     quote in a heading yields clean bookmark text without hyperref warnings.
\pdfstringdefDisableCommands{%
  \def\nwlinkedidentq#1#2{#1}%
  \let\Tt\relax
  \let\nwendquote\relax
}

% Unicode character replacements for symbols not in Latin Modern fonts
% These appear in TUI code and documentation describing terminal output
\usepackage{newunicodechar}
\newunicodechar{✓}{\checkmark}
\newunicodechar{✗}{$\times$}
\newunicodechar{○}{\ensuremath{\circ}}
\newunicodechar{─}{\textemdash}
\newunicodechar{📝}{\texttt{[note]}}
\newunicodechar{🤖}{\texttt{[AI]}}


\title{Scholar CLI: A Tool for Structured Literature Searches}
\author{Daniel Bosk\\Dan-Claude van Thropic\thanks{%
    Wrote this document and the Scholar system.
  } \and
  Ric Glassey\\Claude-Eric Anthropicsson\thanks{%
    Wrote Snowball and Tuxedo.
  }
}
\date{\today}

\begin{document}
\frontmatter
\maketitle

\begin{abstract}
Systematic literature reviews are fundamental to rigorous scholarship, yet
conducting them remains tedious and error-prone.
Researchers must search multiple databases with different interfaces, merge
and deduplicate results, and meticulously document their screening decisions.
Scholar addresses these challenges through a command-line tool that unifies
access to six major academic databases---Semantic Scholar, OpenAlex, DBLP,
Web of Science, IEEE Xplore, and arXiv---while automating the bookkeeping that
systematic reviews demand.

The tool centers on a terminal-based review interface where researchers
navigate papers, read abstracts, record keep/discard decisions with
justifications, and organize findings with themes and notes.
For large result sets, an LLM-assisted classification system learns from
user decisions to suggest classifications for pending papers, dramatically
reducing manual screening effort while maintaining human oversight.
Intelligent caching, automatic deduplication, and PDF management round out
the workflow support.

This document is the complete literate source of Scholar, weaving together
the design rationale, implementation, and tests into a unified narrative.
The code is extracted (``tangled'') from this document using noweb.
\end{abstract}
\vfill
\inputminted[linenos=false]{text}{../LICENSE}

\clearpage
\tableofcontents*

\mainmatter

\input{introduction}

\part{Tutorials}

The chapters in this part are the literate sources of the interactive
tutorials that \texttt{scholar tutorial} runs.
Each tangles to the Markdown that pytorial executes and weaves to a
chapter that explains how the lesson teaches: which aspect of Scholar a
step lets the learner vary, what it holds invariant, and why the steps
come in this order.
That design reasoning is the object of these chapters, so it stands in
the body text; the learner never sees it, because only the Markdown is
shipped.
The tutorials come first because they are the shortest path from a fresh
install to a finished review; the parts that follow explain how the tool
does what these chapters teach.

\input{../src/scholar/tutorials/first-search.tex}
\input{../src/scholar/tutorials/reviewing-in-the-tui.tex}
\input{../src/scholar/tutorials/enrich-verify-export.tex}
\input{../src/scholar/tutorials/llm-synthesis.tex}
\input{../src/scholar/tutorials/scholar-with-an-agent.tex}

\part{Core Functionality}

\documentclass[a4paper,oneside]{memoir}
% Standard preamble for literate programming documents
% Copy this file verbatim to doc/preamble.tex in new projects
%
% This preamble is shared across all literate programming projects.
% The main document (e.g., doc/main.tex) should \input{preamble} after
% \documentclass but before \begin{document}.

% Font setup.  Under pdfLaTeX these are REQUIRED: without [T1]{fontenc} a
% monospace [[code]] quote (or \texttt) cannot select a T1 typewriter shape and
% silently falls back to the proportional OT1 roman font -- i.e. monospace
% "disappears".  iftex guards them so this preamble stays safe when copied into
% a LuaLaTeX/XeLaTeX project (which already use Unicode fonts and need neither).
% Do NOT comment these out.  See latex-writing's references/unicode-and-fonts.md
% ("Symptom 2").
\usepackage{iftex}
\ifpdftex
  \usepackage[utf8]{inputenc}
  \usepackage[T1]{fontenc}
  \usepackage{lmodern}
\fi
\usepackage[swedish,british]{babel}
\usepackage{booktabs}

\usepackage{noweb}
\noweboptions{shift,breakcode,longxref,longchunks}

\usepackage[natbib,style=alphabetic,maxbibnames=99]{biblatex}
\addbibresource{bibliography.bib}

\usepackage[inline]{enumitem}
\setlist[enumerate]{label=(\arabic*)}

\usepackage[all]{foreign}
\renewcommand{\foreignfullfont}{}
\renewcommand{\foreignabbrfont}{}

%\usepackage{newclude}
\usepackage{import}

\usepackage[strict]{csquotes}
\usepackage[single]{acro}

\usepackage{subcaption}

\usepackage[noend]{algpseudocode}
\usepackage{xparse}

\let\email\texttt

%\usepackage[outputdir=ltxobj]{minted}
\usepackage{minted}
\setminted{autogobble}

\usepackage{pythontex}
\setpythontexoutputdir{.}
\setpythontexworkingdir{..}

\usepackage{fancyvrb}

\usepackage{amsmath}
\usepackage{amssymb}
\usepackage{mathtools}
\usepackage{amsthm}
\usepackage{thmtools}
\usepackage[unq]{unique}
\DeclareMathOperator{\powerset}{\mathcal{P}}

\usepackage[binary-units]{siunitx}

\usepackage{pgf}

\usepackage[colorlinks]{hyperref}
\usepackage[capitalize]{cleveref}
\crefalias{question}{item}
\crefname{question}{question}{questions}
\Crefname{question}{Question}{Questions}
\creflabelformat{question}{#2\textup{#1}#3}

% adapted from https://tex.stackexchange.com/a/30726/17418
\newcommand\chnote[1]{%
  \begingroup
  \renewcommand\thefootnote{}\footnote{#1}%
  \addtocounter{footnote}{-1}%
  \endgroup
}

\usepackage[marginparmargin=right]{didactic}

% Make noweb's hyperlinked identifiers safe in headings (chapter/section titles
% that contain a [[code]] quote).  With -index, such a quote weaves into a
% \nwlinkedidentq hyperlink; in a moving argument (the .toc, the running header,
% the PDF bookmark) hyperref then fails fatally (under newer TeX Live with
% errors like "Undefined control sequence \hyper@@link" / "\Hy@reserved@a", or a
% stack overflow from noweb's self-redefining \nwhyperreference).  Two fixes,
% both required, let the project keep [[...]] in headings instead of falling
% back to \texttt{...}:
%
% (1) Replace \nwhyperreference with noweb's robust hyperref form so it is
%     written literally into moving arguments instead of expanding (and
%     misfiring its terminating \global\let) there, yet still links in the body.
\DeclareRobustCommand\nwhyperreference[2]{\hyperlink{noweb.#1}{#2}}
%
% (2) Strip noweb's code-quote macros from PDF bookmark strings, so a [[code]]
%     quote in a heading yields clean bookmark text without hyperref warnings.
\pdfstringdefDisableCommands{%
  \def\nwlinkedidentq#1#2{#1}%
  \let\Tt\relax
  \let\nwendquote\relax
}

% Unicode character replacements for symbols not in Latin Modern fonts
% These appear in TUI code and documentation describing terminal output
\usepackage{newunicodechar}
\newunicodechar{✓}{\checkmark}
\newunicodechar{✗}{$\times$}
\newunicodechar{○}{\ensuremath{\circ}}
\newunicodechar{─}{\textemdash}
\newunicodechar{📝}{\texttt{[note]}}
\newunicodechar{🤖}{\texttt{[AI]}}


\title{Scholar CLI: A Tool for Structured Literature Searches}
\author{Daniel Bosk\\Dan-Claude van Thropic\thanks{%
    Wrote this document and the Scholar system.
  } \and
  Ric Glassey\\Claude-Eric Anthropicsson\thanks{%
    Wrote Snowball and Tuxedo.
  }
}
\date{\today}

\begin{document}
\frontmatter
\maketitle

\begin{abstract}
Systematic literature reviews are fundamental to rigorous scholarship, yet
conducting them remains tedious and error-prone.
Researchers must search multiple databases with different interfaces, merge
and deduplicate results, and meticulously document their screening decisions.
Scholar addresses these challenges through a command-line tool that unifies
access to six major academic databases---Semantic Scholar, OpenAlex, DBLP,
Web of Science, IEEE Xplore, and arXiv---while automating the bookkeeping that
systematic reviews demand.

The tool centers on a terminal-based review interface where researchers
navigate papers, read abstracts, record keep/discard decisions with
justifications, and organize findings with themes and notes.
For large result sets, an LLM-assisted classification system learns from
user decisions to suggest classifications for pending papers, dramatically
reducing manual screening effort while maintaining human oversight.
Intelligent caching, automatic deduplication, and PDF management round out
the workflow support.

This document is the complete literate source of Scholar, weaving together
the design rationale, implementation, and tests into a unified narrative.
The code is extracted (``tangled'') from this document using noweb.
\end{abstract}
\vfill
\inputminted[linenos=false]{text}{../LICENSE}

\clearpage
\tableofcontents*

\mainmatter

\input{introduction}

\part{Tutorials}

The chapters in this part are the literate sources of the interactive
tutorials that \texttt{scholar tutorial} runs.
Each tangles to the Markdown that pytorial executes and weaves to a
chapter that explains how the lesson teaches: which aspect of Scholar a
step lets the learner vary, what it holds invariant, and why the steps
come in this order.
That design reasoning is the object of these chapters, so it stands in
the body text; the learner never sees it, because only the Markdown is
shipped.
The tutorials come first because they are the shortest path from a fresh
install to a finished review; the parts that follow explain how the tool
does what these chapters teach.

\input{../src/scholar/tutorials/first-search.tex}
\input{../src/scholar/tutorials/reviewing-in-the-tui.tex}
\input{../src/scholar/tutorials/enrich-verify-export.tex}
\input{../src/scholar/tutorials/llm-synthesis.tex}
\input{../src/scholar/tutorials/scholar-with-an-agent.tex}

\part{Core Functionality}

\documentclass[a4paper,oneside]{memoir}
\input{preamble}

\title{Scholar CLI: A Tool for Structured Literature Searches}
\author{Daniel Bosk\\Dan-Claude van Thropic\thanks{%
    Wrote this document and the Scholar system.
  } \and
  Ric Glassey\\Claude-Eric Anthropicsson\thanks{%
    Wrote Snowball and Tuxedo.
  }
}
\date{\today}

\begin{document}
\frontmatter
\maketitle

\begin{abstract}
Systematic literature reviews are fundamental to rigorous scholarship, yet
conducting them remains tedious and error-prone.
Researchers must search multiple databases with different interfaces, merge
and deduplicate results, and meticulously document their screening decisions.
Scholar addresses these challenges through a command-line tool that unifies
access to six major academic databases---Semantic Scholar, OpenAlex, DBLP,
Web of Science, IEEE Xplore, and arXiv---while automating the bookkeeping that
systematic reviews demand.

The tool centers on a terminal-based review interface where researchers
navigate papers, read abstracts, record keep/discard decisions with
justifications, and organize findings with themes and notes.
For large result sets, an LLM-assisted classification system learns from
user decisions to suggest classifications for pending papers, dramatically
reducing manual screening effort while maintaining human oversight.
Intelligent caching, automatic deduplication, and PDF management round out
the workflow support.

This document is the complete literate source of Scholar, weaving together
the design rationale, implementation, and tests into a unified narrative.
The code is extracted (``tangled'') from this document using noweb.
\end{abstract}
\vfill
\inputminted[linenos=false]{text}{../LICENSE}

\clearpage
\tableofcontents*

\mainmatter

\input{introduction}

\part{Tutorials}

The chapters in this part are the literate sources of the interactive
tutorials that \texttt{scholar tutorial} runs.
Each tangles to the Markdown that pytorial executes and weaves to a
chapter that explains how the lesson teaches: which aspect of Scholar a
step lets the learner vary, what it holds invariant, and why the steps
come in this order.
That design reasoning is the object of these chapters, so it stands in
the body text; the learner never sees it, because only the Markdown is
shipped.
The tutorials come first because they are the shortest path from a fresh
install to a finished review; the parts that follow explain how the tool
does what these chapters teach.

\input{../src/scholar/tutorials/first-search.tex}
\input{../src/scholar/tutorials/reviewing-in-the-tui.tex}
\input{../src/scholar/tutorials/enrich-verify-export.tex}
\input{../src/scholar/tutorials/llm-synthesis.tex}
\input{../src/scholar/tutorials/scholar-with-an-agent.tex}

\part{Core Functionality}

\input{../src/scholar/scholar.tex}
\input{../src/scholar/utils.tex}
\input{../src/scholar/cache.tex}
\input{../src/scholar/ratelimit.tex}

\input{../src/scholar/providers.tex}
\input{../src/scholar/enrich.tex}
\input{../src/scholar/crossref.tex}
\input{../src/scholar/unpaywall.tex}
\input{../src/scholar/preprint.tex}
\input{../src/scholar/pdf.tex}

\part{Review System}

\input{../src/scholar/notes.tex}
\input{../src/scholar/review.tex}
\input{../src/scholar/llm_review.tex}

\part{User Interface}

\input{../src/scholar/cli.tex}
\input{../src/scholar/tui.tex}

\backmatter

\end{document}

\input{../src/scholar/utils.tex}
\input{../src/scholar/cache.tex}
\input{../src/scholar/ratelimit.tex}

\input{../src/scholar/providers.tex}
\input{../src/scholar/enrich.tex}
\input{../src/scholar/crossref.tex}
\input{../src/scholar/unpaywall.tex}
\input{../src/scholar/preprint.tex}
\input{../src/scholar/pdf.tex}

\part{Review System}

\input{../src/scholar/notes.tex}
\input{../src/scholar/review.tex}
\input{../src/scholar/llm_review.tex}

\part{User Interface}

\input{../src/scholar/cli.tex}
\input{../src/scholar/tui.tex}

\backmatter

\end{document}

\input{../src/scholar/utils.tex}
\input{../src/scholar/cache.tex}
\input{../src/scholar/ratelimit.tex}

\input{../src/scholar/providers.tex}
\input{../src/scholar/enrich.tex}
\input{../src/scholar/crossref.tex}
\input{../src/scholar/unpaywall.tex}
\input{../src/scholar/preprint.tex}
\input{../src/scholar/pdf.tex}

\part{Review System}

\input{../src/scholar/notes.tex}
\input{../src/scholar/review.tex}
\input{../src/scholar/llm_review.tex}

\part{User Interface}

\input{../src/scholar/cli.tex}
\input{../src/scholar/tui.tex}

\backmatter

\end{document}

\input{../src/scholar/utils.tex}
\input{../src/scholar/cache.tex}
\input{../src/scholar/ratelimit.tex}

\input{../src/scholar/providers.tex}
\input{../src/scholar/enrich.tex}
\input{../src/scholar/crossref.tex}
\input{../src/scholar/unpaywall.tex}
\input{../src/scholar/preprint.tex}
\input{../src/scholar/pdf.tex}

\part{Review System}

\input{../src/scholar/notes.tex}
\input{../src/scholar/review.tex}
\input{../src/scholar/llm_review.tex}

\part{User Interface}

\input{../src/scholar/cli.tex}
\input{../src/scholar/tui.tex}

\backmatter

\end{document}
